\section{Task 2: Chatbot Models}
\label{sec:task2}

Task 2 trains two sequence-to-sequence models to answer factual questions. Both
share the same vocabulary and dataset split.

\subsection{LSTM Encoder-Decoder with Bahdanau Attention}

The LSTM chatbot (\texttt{models/lstm\_chatbot.py}) follows the classic
encoder-decoder paradigm with an attention mechanism.

\subsubsection*{Encoder}

A \textbf{bidirectional} 2-layer LSTM reads the input question token by token.
Bidirectionality means the encoder processes each word in both left-to-right and
right-to-left order, so each output position has context from both sides of the
sentence. The forward and backward hidden states are concatenated and projected
to the decoder's hidden dimension via a linear layer.

\subsubsection*{Attention}

At each decoding step, Bahdanau attention assigns a weight to every encoder
output, indicating how relevant that input token is for generating the current
output word. These weights produce a context vector — a weighted summary of the
input — which is passed to the decoder alongside the previous predicted word.
This allows the model to focus on different parts of the question at each step.

\subsubsection*{Decoder}

A unidirectional LSTM decoder generates the answer one word at a time. At each
step it receives the previous predicted word and the attention context vector,
then predicts the next word from the vocabulary.

\subsubsection*{Teacher Forcing}

During training, the correct previous word (rather than the model's own
prediction) is fed to the decoder. This technique, called teacher forcing,
provides a strong training signal early on. The ratio decreases from 1.0 to 0.5
during training, gradually encouraging the model to rely on its own outputs.

\textbf{Architecture summary:}
\begin{itemize}
    \item Embedding dimension: 128
    \item Hidden dimension: 256 (bidirectional encoder: 512 total)
    \item Total parameters: 6,019,947
\end{itemize}

\subsection{Transformer Encoder-Decoder}

The Transformer chatbot (\texttt{models/transformer\_chatbot.py}) implements the
full encoder-decoder Transformer \cite{vaswani2017} from scratch.

\subsubsection*{Positional Encoding}

Since the Transformer processes all tokens in parallel rather than sequentially,
it has no built-in sense of word order. Sinusoidal positional encodings are added
to the input embeddings to provide this information — each position gets a unique
pattern of sine and cosine values that the model can learn to interpret.

\subsubsection*{Encoder and Decoder Stacks}

Each encoder layer applies multi-head self-attention (every word attends to every
other word) followed by a feed-forward network, with residual connections and
layer normalisation around each sub-layer. The decoder has the same structure but
adds a cross-attention layer that attends over the encoder outputs, connecting the
question representation to the answer being generated.

\subsubsection*{Learning Rate Schedule}

A warmup schedule is used: the learning rate increases gradually for the first
few hundred steps, then decays. This prevents the model from making large
destructive weight updates at the start of training when its parameters are still
random.

\textbf{Architecture summary:}
\begin{itemize}
    \item Model dimension: 128
    \item Feed-forward dimension: 512
    \item Attention heads: 4
    \item Encoder layers: 3 \quad Decoder layers: 3
    \item Total parameters: 1,971,819
\end{itemize}

\subsection{Training Configuration}

Both chatbots are trained for \textbf{30 epochs} on the same 2,000+ QA pairs with:
\begin{itemize}
    \item Optimiser: Adam (LSTM), Adam with warmup schedule (Transformer)
    \item Gradient clipping: max norm $= 1.0$
    \item Loss: cross-entropy (padding tokens ignored)
\end{itemize}
